% ============================================================
% BIBLIOGRAPHIC ORGANIZATION
% Viability Philosophy System
% Book I: Ontology of Viability
% ============================================================

\cleardoublepage

\section*{Bibliographic Organization}

\phantomsection
\addcontentsline{toc}{section}{Bibliographic Organization}

% ------------------------------------------------------------
% 1. PURPOSE AND SCOPE
% ------------------------------------------------------------

\subsection*{1. Purpose and Scope}

Bibliographic organization establishes
the principles governing the systematic
recording, arrangement, presentation,
and maintenance of external sources
used within the book.

Its purpose is to ensure that external
works cited in the text can be identified,
distinguished, retrieved, and verified.

Bibliographic organization provides
the documentary structure corresponding
to the external source references
discussed in the preceding document.

It does not establish a new internal
identification system for external
works. Rather, it organizes their
bibliographic descriptions according
to consistent and appropriate
conventions.

The bibliographic organization of
the book must support the following
requirements:

\begin{itemize}
    \item accurate identification
    of external works;
    \item consistency between
    in-text citations and
    bibliographic records;
    \item sufficient information
    for source retrieval;
    \item distinction between
    different editions and
    versions of a work;
    \item systematic arrangement
    of bibliographic entries;
    \item reliable maintenance
    during manuscript revision;
    \item transparency concerning
    the external materials used
    in the philosophical discussion.
\end{itemize}

This document establishes the
organizational principles of
bibliographic records.

The precise citation style,
typographical conventions,
and technical implementation
must be determined by the
bibliographic style adopted
for the book.

% ------------------------------------------------------------
% 2. BIBLIOGRAPHIC RECORDS
% ------------------------------------------------------------

\subsection*{2. Bibliographic Records}

A bibliographic record represents
an external work or document
through the information necessary
for its identification and retrieval.

Each record must contain sufficient
information to distinguish the
source from other works and to
identify the version relevant
to the discussion.

Depending on the source type,
a bibliographic record may include:

\begin{itemize}
    \item author or responsible creator;
    \item title and subtitle;
    \item editor, translator,
    or other relevant contributor;
    \item edition or version;
    \item publication date;
    \item publisher or issuing institution;
    \item journal, collection,
    or larger publication;
    \item volume and issue information;
    \item page range, where applicable;
    \item persistent identifier,
    where available;
    \item stable retrieval information,
    where appropriate.
\end{itemize}

The required elements depend
on the nature of the source
and the bibliographic conventions
adopted for the book.

A record must not contain
invented or unverified
bibliographic information.

Where an element is unavailable,
its treatment must follow the
adopted bibliographic conventions.

% ------------------------------------------------------------
% 3. DISTINCTION BETWEEN WORKS,
% EDITIONS, AND VERSIONS
% ------------------------------------------------------------

\subsection*{3. Distinction Between Works,
Editions, and Versions}

A bibliographic organization
must preserve the distinction
between a work and the particular
edition or version consulted.

A work may be available in
multiple editions, translations,
revisions, or digital versions.

These forms may differ in:

\begin{itemize}
    \item textual content;
    \item pagination;
    \item organization;
    \item editorial interventions;
    \item translation choices;
    \item publication information.
\end{itemize}

When such differences are
relevant to the identification
or verification of a cited passage,
the bibliographic record must
identify the version actually
consulted.

Different editions must not
be treated as interchangeable
when their differences affect
the reference or interpretation.

Where multiple versions of
the same work are used,
their bibliographic records
must preserve the distinctions
necessary for accurate citation.

% ------------------------------------------------------------
% 4. TYPES OF BIBLIOGRAPHIC
% MATERIAL
% ------------------------------------------------------------

\subsection*{4. Types of Bibliographic Material}

The bibliography may include
different kinds of external
materials used in the book.

Depending on the sources
actually consulted, these
may include:

\begin{itemize}
    \item books and monographs;
    \item chapters in edited volumes;
    \item journal articles;
    \item conference papers
    and proceedings;
    \item institutional reports;
    \item archival documents;
    \item primary historical texts;
    \item translated works;
    \item digital publications;
    \item online documents;
    \item audiovisual materials;
    \item other identifiable
    documentary sources.
\end{itemize}

The bibliographic description
must be appropriate to the
type of material being recorded.

Different source types may
require different bibliographic
elements and different
arrangements of information.

The existence of a source
category in this list does
not imply that the book
must contain sources of
that category.

Only materials actually
relevant to the manuscript
should be included.

% ------------------------------------------------------------
% 5. RELATION BETWEEN
% IN-TEXT CITATIONS AND
% BIBLIOGRAPHIC RECORDS
% ------------------------------------------------------------

\subsection*{5. Relation Between In-Text
Citations and Bibliographic Records}

In-text citations and
bibliographic records
perform complementary
functions.

An in-text citation identifies
the external source relevant
to a particular passage
or discussion.

A bibliographic record
provides the information
necessary to identify
and retrieve that source.

The two must remain
consistent and mutually
traceable.

Where an in-text citation
refers to a bibliographic
record, the corresponding
record must be identifiable
without ambiguity.

Likewise, bibliographic
records must not contain
unexplained inconsistencies
with the citations used
in the text.

The relationship between
citations and records must
be reviewed whenever
sources are added, removed,
or revised.

The exact form of this
relationship depends on
the citation style adopted
for the book.

% ------------------------------------------------------------
% 6. ORGANIZATION AND
% ARRANGEMENT OF RECORDS
% ------------------------------------------------------------

\subsection*{6. Organization and Arrangement
of Records}

Bibliographic records must
be arranged systematically
so that readers can locate
the required sources.

The arrangement must be
consistent throughout
the bibliography.

Possible organizational
principles include:

\begin{itemize}
    \item alphabetical arrangement
    by author or responsible
    institution;
    \item chronological arrangement;
    \item separation by source
    category;
    \item another systematic
    arrangement permitted
    by the adopted bibliographic
    style.
\end{itemize}

The final arrangement
must be selected as part
of the book's bibliographic
conventions.

Different organizational
principles must not be
combined arbitrarily.

If the bibliography is
divided into categories
or sections, the basis
for that division must
be clear and consistently
applied.

The organization must
facilitate retrieval without
creating unnecessary
duplication or ambiguity.

% ------------------------------------------------------------
% 7. AUTHOR AND RESPONSIBLE
% ENTITY IDENTIFICATION
% ------------------------------------------------------------

\subsection*{7. Author and Responsible
Entity Identification}

Authors and responsible
institutions must be
represented consistently
in bibliographic records.

The form of a person's
name or an institution's
name must permit reliable
identification of the
corresponding source.

Where a source has multiple
authors, editors, translators,
or other contributors,
their roles must be
represented according
to the adopted bibliographic
conventions.

An institution acting
as the responsible creator
or publisher must be
identified appropriately.

Differences in spelling,
transliteration, name order,
or institutional naming
must be handled consistently.

A bibliographic record
must not assign authorship
or responsibility without
adequate support from
the source itself or
reliable bibliographic
information.

% ------------------------------------------------------------
% 8. TITLES AND PUBLICATION
% INFORMATION
% ------------------------------------------------------------

\subsection*{8. Titles and Publication Information}

Titles must be transcribed
or represented accurately
according to the adopted
bibliographic conventions.

Where a source has a
subtitle, translated title,
alternative title, or
other identifying title
information, its treatment
must preserve the
distinctions relevant
to identification.

Publication information
must correspond to the
edition or version
actually consulted.

The bibliographic record
must not combine publication
details from different
editions or versions
without an explicit
and justified reason.

Where a source has
incomplete publication
information, the missing
elements must be handled
according to the adopted
style rather than completed
through speculation.

% ------------------------------------------------------------
% 9. TRANSLATIONS AND
% TRANSLATORS
% ------------------------------------------------------------

\subsection*{9. Translations and Translators}

Translated works require
particular attention because
the identity of the original
work and the identity
of the consulted translation
are related but distinct.

Where a published translation
is consulted, the bibliographic
record must identify the
translation and its relevant
publication information.

The translator should be
identified where required
by the adopted bibliographic
conventions or where the
translator's contribution
is relevant to the discussion.

When the author of the
present book provides
a translation directly
from an original-language
source, that fact must
be distinguished from
the use of a published
translation where relevant.

The bibliography must
not obscure which textual
version supports a quotation
or interpretation.

% ------------------------------------------------------------
% 10. MULTIPLE WORKS BY
% THE SAME AUTHOR
% ------------------------------------------------------------

\subsection*{10. Multiple Works by
the Same Author}

When multiple works by
the same author are included,
their bibliographic records
must remain distinguishable.

The arrangement of these
records must follow the
selected bibliographic
convention.

Where works share the
same author and publication
date, the adopted style
may require additional
distinguishing elements.

Such distinctions must
be applied consistently
both in the bibliography
and in the corresponding
in-text citations.

The purpose of these
distinctions is to prevent
ambiguity when the reader
moves between a citation
and its bibliographic record.

% ------------------------------------------------------------
% 11. DIGITAL SOURCES AND
% PERSISTENT IDENTIFIERS
% ------------------------------------------------------------

\subsection*{11. Digital Sources and
Persistent Identifiers}

Digital sources must be
recorded in a manner
that supports their
identification and retrieval.

Where available and
appropriate, a bibliographic
record may include:

\begin{itemize}
    \item a Digital Object
    Identifier (DOI);
    \item another persistent
    identifier;
    \item a stable URL;
    \item version or release
    information;
    \item publication or
    update information;
    \item an access date,
    where required.
\end{itemize}

A persistent identifier
should be preferred when
it is available and suitable
for identifying the cited
material.

A URL must be represented
accurately and must not
be treated as a substitute
for the other identifying
information required
by the source type.

Where digital material
is subject to change,
the relevant version
or retrieval context
must be identified
when necessary.

% ------------------------------------------------------------
% 12. PRIMARY AND SECONDARY
% SOURCES
% ------------------------------------------------------------

\subsection*{12. Primary and Secondary Sources}

The bibliography may contain
both primary and secondary
sources, depending on
the nature of the inquiry.

The distinction between
primary and secondary
sources concerns their
function and relation
to the subject under
investigation.

A primary source may
constitute the original
material being examined.

A secondary source may
provide interpretation,
analysis, historical
context, or critical
discussion of another
work or subject.

These categories must
not be confused with
the distinction between
direct and indirect
consultation.

A source may be primary
in relation to the subject
of inquiry while still
being consulted through
a particular edition
or translation.

Likewise, a secondary
source may be consulted
directly.

The classification of
sources must therefore
reflect their actual
function within the
research and argument.

Where the bibliography
uses separate sections
for primary and secondary
materials, the distinction
must be applied consistently.

% ------------------------------------------------------------
% 13. DUPLICATE RECORDS
% ------------------------------------------------------------

\subsection*{13. Duplicate Records}

Bibliographic records
must be reviewed for
unnecessary duplication.

The same edition or
version of a work should
not appear as multiple
independent records
merely because it is
cited in different
chapters or sections.

Different editions,
translations, or versions
must remain distinguishable
when they represent
different consulted
materials.

The removal of duplicate
records must not eliminate
bibliographic distinctions
that are necessary for
accurate citation.

Where a work appears
in multiple forms,
the bibliographic organization
must preserve the identity
of each form actually
used in the manuscript.

% ------------------------------------------------------------
% 14. COMPLETENESS AND
% VERIFICATION
% ------------------------------------------------------------

\subsection*{14. Completeness and Verification}

Bibliographic records
must be checked for
accuracy and completeness
before the manuscript
is finalized.

Verification should
address the following:

\begin{itemize}
    \item accuracy of author
    and contributor names;
    \item accuracy of titles;
    \item correctness of
    edition and version
    information;
    \item publication details;
    \item consistency of
    dates and identifiers;
    \item accuracy of URLs
    and persistent identifiers;
    \item correspondence
    between citations
    and bibliographic
    records;
    \item absence of
    unnecessary duplicates.
\end{itemize}

A record must not be
considered verified
merely because it has
been formatted correctly.

Bibliographic accuracy
depends on the correctness
of the underlying
source information.

Where verification
is incomplete, uncertainty
must not be concealed
through fabricated
or assumed details.

% ------------------------------------------------------------
% 15. BIBLIOGRAPHIC RECORDS
% AND INTERNAL IDENTIFICATION
% ------------------------------------------------------------

\subsection*{15. Bibliographic Records
and Internal Identification}

Bibliographic records
must remain separate
from the internal
identification systems
of the Viability Philosophy
System.

An external source
does not receive an
internal structural
identifier merely
because it is cited
in a particular location.

Likewise, it does not
receive an internal
logical entity identifier
merely because its
concepts, arguments,
or author are discussed
within the book.

The internal structural
identifier identifies
a location within the
book.

The internal logical
entity identifier
identifies an entity
belonging to the
system's internal
logical organization.

The bibliographic record
identifies an external
work or document.

These functions must
not be conflated.

The relationship between
an internal entity and
an external source
must be expressed
through the relevant
argument, textual
references, and
bibliographic citations.

% ------------------------------------------------------------
% 16. BIBLIOGRAPHY AND
% SOURCE USE
% ------------------------------------------------------------

\subsection*{16. Bibliography and Source Use}

The bibliography must
reflect the actual use
of external materials
within the manuscript
according to the
bibliographic convention
adopted for the book.

The inclusion of a
source must have
a legitimate documentary
or scholarly function.

A bibliographic record
must not be added
merely to create an
appearance of extensive
research.

Likewise, sources
that support relevant
claims, quotations,
or interpretations
must not be omitted
from the appropriate
references.

The precise relationship
between works cited
in the text and other
consulted materials
must be determined
by the selected
bibliographic style.

The book must apply
that distinction
consistently.

% ------------------------------------------------------------
% 17. REVISION AND
% MAINTENANCE
% ------------------------------------------------------------

\subsection*{17. Revision and Maintenance}

Bibliographic records
must be maintained
throughout the development
and revision of the manuscript.

The following events
require bibliographic
review:

\begin{itemize}
    \item the addition
    of a new external source;
    \item the removal
    of an existing source;
    \item the replacement
    of an edition or
    translation;
    \item the correction
    of publication information;
    \item the revision
    of an in-text citation;
    \item the addition
    or removal of a
    quotation or paraphrase;
    \item the restructuring
    of bibliographic
    sections;
    \item the identification
    of duplicate or
    inconsistent records.
\end{itemize}

Whenever the manuscript
is revised, the
bibliographic records
and their corresponding
citations must be
reviewed for consistency.

Bibliographic maintenance
is part of the intellectual
and documentary revision
of the book, not merely
a final formatting task.

% ------------------------------------------------------------
% 18. BIBLIOGRAPHIC STYLE
% AND TECHNICAL IMPLEMENTATION
% ------------------------------------------------------------

\subsection*{18. Bibliographic Style and
Technical Implementation}

The bibliographic principles
defined in this document
must be implemented
through a consistent
citation and bibliography
style.

The adopted style
determines matters such as:

\begin{itemize}
    \item the form of
    in-text citations;
    \item the ordering
    of bibliographic
    elements;
    \item punctuation
    and typographical
    conventions;
    \item treatment of
    multiple authors;
    \item treatment of
    editions and translations;
    \item ordering of
    bibliographic entries;
    \item formatting of
    digital sources;
    \item distinction
    between different
    bibliographic sections.
\end{itemize}

The technical implementation
may use an appropriate
bibliographic management
system or typesetting
workflow.

However, the technical
tool must serve the
bibliographic principles
rather than determine
the conceptual distinctions
independently.

The selected implementation
must support consistency,
traceability, and reliable
maintenance throughout
the book.

The specific citation
style and technical
implementation remain
subject to an explicit
project-level decision.

% ------------------------------------------------------------
% 19. RELATION TO THE OTHER
% REFERENCING DOCUMENTS
% ------------------------------------------------------------

\subsection*{19. Relation to the Other
Referencing Documents}

Bibliographic organization
completes the framework
established by the
other documents in
this guide.

Its relation to those
documents may be
summarized as follows:

\begin{itemize}
    \item \textit{Structural Identifiers}
    establish the identification
    of internal textual locations.

    \item \textit{Logical Entity
    Identifiers} establish
    the identification
    of internal logical
    entities.

    \item \textit{Internal
    Cross-References}
    establish references
    between locations
    and entities within
    the system.

    \item \textit{Entity
    Location References}
    connect logical
    entities with their
    textual occurrences.

    \item \textit{Logical
    Dependencies} describe
    relations of logical
    dependence between
    internal entities.

    \item \textit{External
    Source References}
    establish the principles
    for citing and locating
    external materials.

    \item \textit{Bibliographic
    Organization} establishes
    the systematic
    organization and
    maintenance of
    the corresponding
    bibliographic records.
\end{itemize}

Together, these documents
provide a framework
for distinguishing
internal structure,
internal logical identity,
internal relations,
external references,
and bibliographic
documentation.

% ------------------------------------------------------------
% 20. SUMMARY
% ------------------------------------------------------------

\subsection*{20. Summary}

Bibliographic organization
provides the systematic
framework for recording,
arranging, presenting,
and maintaining the
external sources used
within the book.

Each bibliographic
record must identify
its source accurately
and provide sufficient
information for
retrieval and verification.

The organization
must preserve distinctions
between works, editions,
translations, and
versions where
those distinctions
are relevant.

In-text citations
and bibliographic
records must remain
consistent and
mutually traceable.

The arrangement
of records, the
treatment of different
source types, and
the technical form
of citations must
follow a coherent
bibliographic convention.

Bibliographic records
must remain distinct
from the internal
structural and
logical identification
systems of the
Viability Philosophy
System.

The resulting framework
connects the book's
internal philosophical
discussion to external
materials in a manner
that is accurate,
transparent, systematic,
and verifiable.

\cleardoublepage

% ============================================================
% END OF BIBLIOGRAPHIC ORGANIZATION
% ============================================================